\section*{Hoofdstuk \ref{ch:b2:blood-circulation} --- Bloed en de bloedsomloop}

\begin{solution}{exo:b2:blood-circulation:1}
\omterm{def:b2:blood-circulation:blood}{Plasma} \qty{55}{\%} (water, \qty{70}{g/L} eiwit, zouten, voedingsstoffen,
gassen, hormonen); cellen \qty{45}{\%}. Rode bloedcellen $5\times
10^{12}$/L, schijfjes van \qty{7}{\micro m}, zuurstoftransport; witte
bloedcellen $7\times 10^{9}$/L, \qtyrange{10}{15}{\micro m}, immuniteit;
\omterm{def:b2:blood-circulation:blood}{bloedplaatjes} $3\times 10^{11}$/L, fragmenten van \qty{2}{\micro m},
hemostase.
\end{solution}

\begin{solution}{exo:b2:blood-circulation:2}
Vis: hart $\to$ kieuwen $\to$ lichaam $\to$ hart, één kringloop, waarbij
het lichaam \omterm{def:b2:blood-circulation:blood}{bloed} krijgt met de lage druk die na de kieuwen overblijft.
Zoogdier: rechterhart $\to$ longen (\qty{25}{mmHg}) $\to$ linkerhart $\to$
lichaam (\qty{120}{mmHg}) $\to$ rechterhart. De dubbele kringloop geeft het
lichaam een hoge druk zonder de longen eraan bloot te stellen, en laat de
twee kringlopen afzonderlijk regelen.
\end{solution}

\begin{solution}{exo:b2:blood-circulation:3}
\omterm{def:b2:blood-circulation:vessels}{Slagader}: dikke wand met elastische lagen en spier --- geleidt \omterm{def:b2:blood-circulation:blood}{bloed} onder
hoge druk en vlakt de pols af. \omterm{def:b2:blood-circulation:vessels}{Arteriole}: wand vooral uit glad spierweefsel
--- bepaalt de weerstand en verdeelt de stroom. \omterm{def:b2:blood-circulation:vessels}{Haarvat}: één endotheelcel
dik --- uitwisseling. \omterm{def:b2:blood-circulation:vessels}{Ader}: dunne, rekbare wand met kleppen --- voert
\omterm{def:b2:blood-circulation:blood}{bloed} onder lage druk terug en bevat het grootste deel ervan.
\end{solution}

\begin{solution}{exo:b2:blood-circulation:4}
$70\times 72 = \qty{5}{L/min}$, \qty{7200}{L} per dag, tegenover
\qty{5}{L} \omterm{def:b2:blood-circulation:blood}{bloed}: het volume gaat \num{1440} keer per dag door het hart.
Geen enkel orgaan zou zeven ton \omterm{def:b2:blood-circulation:blood}{bloed} per dag kunnen aanmaken of verbruiken;
hetzelfde \omterm{def:b2:blood-circulation:blood}{bloed} moet terugkeren --- het circuleert.
\end{solution}

\begin{solution}{exo:b2:blood-circulation:5}
$\Delta P = 8\eta LQ/\pi r^{4} = 8\times 3\times 10^{-3}\times
0.4\times 8.3\times 10^{-5}/(\pi\times 2.44\times 10^{-8}) =
\qty{10}{Pa}$, \qty{0.08}{mmHg}: de aorta kost niets; de druk wordt in de
\omterm{def:b2:blood-circulation:vessels}{arteriolen} verbruikt.
\end{solution}

\begin{solution}{exo:b2:blood-circulation:6}
$(15/12)^{4} = 2.4$: weerstand 2.4 keer zo groot, stroom gedaald tot
\qty{41}{\%}. $(15/20)^{4} = 0.32$: weerstand gedaald tot een derde, stroom
3.2 keer zo groot.
\end{solution}

\begin{solution}{exo:b2:blood-circulation:7}
Aorta $83/3 = \qty{28}{cm/s}$; \omterm{def:b2:blood-circulation:vessels}{haarvaten} $83/3000 =
\qty{0.028}{cm/s}$; doorlooptijd $0.1/0.028 = \qty{3.6}{s}$. Bij inspanning
stijgt het debiet vijfvoudig en het open haarvatoppervlak misschien
drievoudig, zodat de doorlooptijd daalt tot ongeveer \qty{2}{s} --- nog
steeds genoeg om de zuurstof te laten vertrekken.
\end{solution}

\begin{solution}{exo:b2:blood-circulation:8}
Arterieel uiteinde $35 - 25 = \qty{+10}{mmHg}$ (filtratie); veneus uiteinde $15
- 25 = \qty{-10}{mmHg}$ (resorptie). Met $P_c = 28$ aan het veneuze
uiteinde is het netto $+3$: over de hele lengte filtreert vloeistof en wordt
er niets geresorbeerd --- oedeem, de gezwollen enkels van hartfalen.
\end{solution}

\begin{solution}{exo:b2:blood-circulation:9}
$40/66000 = \qty{0.61}{mol/m^3}$; $\pi = 0.61\times 8.314\times 310 =
\qty{1570}{Pa} = \qty{11.8}{mmHg}$; bij de helft van de albumine
\qty{5.9}{mmHg}. Natrium passeert de haarvatwand vrij, zodat zijn
concentratie aan beide kanten gelijk is en het er geen osmotische druk over
uitoefent.
\end{solution}

\begin{solution}{exo:b2:blood-circulation:10}
$RC = \qty{2}{s}$: $120e^{-0.4} = \qty{80}{mmHg}$. $C = 1$: $RC = 1$,
$120e^{-0.8} = \qty{54}{mmHg}$ --- en hetzelfde slagvolume verhoogt de
systolische druk in een stijve aorta sterker. Een grote \omterm{thm:b2:blood-circulation:windkessel}{polsdruk} betekent
een hogere piek waartegen het hart moet uitpompen, meer arbeid en meer
zuurstof voor hetzelfde debiet.
\end{solution}

\begin{solution}{exo:b2:blood-circulation:11}
Rust: $90/83 = \qty{1.08}{mmHg.s/mL}$; inspanning: $105/417 =
\qty{0.25}{mmHg.s/mL}$, een viervoudige daling. Omdat $R \propto r^{-4}$,
is $r$ gestegen met een factor $4^{1/4} = 1.41$: de \omterm{def:b2:blood-circulation:vessels}{arteriolen} zijn
gemiddeld \qty{41}{\%} wijder geworden.
\end{solution}

\begin{solution}{exo:b2:blood-circulation:12}
De continuïteit legt op elk niveau de snelheid vast via de totale doorsnede,
en de boom is zo gebouwd dat die doorsnede in de \omterm{def:b2:blood-circulation:vessels}{haarvaten} duizend keer die
van de aorta is en in de \omterm{def:b2:blood-circulation:vessels}{aders} weer klein; de \omterm{thm:b2:blood-circulation:poiseuille}{wet van Poiseuille} legt de
weerstand in de \omterm{def:b2:blood-circulation:vessels}{arteriolen}, waar spier hem kan veranderen, en laat de wijde
vaten vrijwel zonder verlies. Een \omterm{def:b2:blood-circulation:circuits}{open bloedsomloop} heeft geen \omterm{def:b2:blood-circulation:vessels}{haarvaten} om
het \omterm{def:b2:blood-circulation:blood}{bloed} op een welbepaalde plaats te vertragen, geen vaten om een
weerstand in te stellen, en geen manier om de stroom naar één orgaan te
leiden --- hij kan alleen roeren.
\end{solution}

\begin{solution}{pb:b2:blood-circulation:1}
\textbf{1.} $70\times 72 = \qty{5.0}{L/min}$; \qty{7300}{L} per dag.
\textbf{2.} $7300/5 \approx 1450$ keer.
\textbf{3.} Een kwart van \qty{57}{g} bij 72 slagen per minuut: \qty{14}{g}
$\times 4320 = \qty{62}{kg}$ per uur --- ongeveer het gewicht van een man.
\textbf{4.} Geen voedsel zou elk uur het gewicht van een man aan \omterm{def:b2:blood-circulation:blood}{bloed} kunnen
leveren, en geen weefsel het kunnen verbruiken; de enige uitweg is dat
hetzelfde \omterm{def:b2:blood-circulation:blood}{bloed} naar het hart terugkeert.
\textbf{5.} Een koord om de arm doet de \omterm{def:b2:blood-circulation:vessels}{aders} eronder opzwellen, niet
erboven, dus beweegt het \omterm{def:b2:blood-circulation:blood}{bloed} in de \omterm{def:b2:blood-circulation:vessels}{aders} naar het hart; wanneer men het
\omterm{def:b2:blood-circulation:blood}{bloed} in een \omterm{def:b2:blood-circulation:vessels}{ader} naar de hand duwt, stopt het bij de kleppen en kan het
niet worden teruggeduwd --- de kleppen laten alleen stroming naar het hart
toe.
\textbf{6.} De \omterm{def:b2:blood-circulation:vessels}{haarvaten} die \omterm{def:b2:blood-circulation:vessels}{slagaders} met \omterm{def:b2:blood-circulation:vessels}{aders} verbinden; Malpighi zag ze
in 1661 in de long van een kikker.
\textbf{7.} Oppervlak $\pi\times 1.25^{2} = \qty{4.9}{cm^2}$; $83/4.9 =
\qty{17}{cm/s}$.
\textbf{8.} $8\times 3\times 10^{-3}\times 0.4\times 8.3\times
10^{-5}/(\pi\times 2.44\times 10^{-8}) = \qty{10}{Pa}$, \qty{0.08}{mmHg}.
\textbf{9.} $R_{1} = 8\times 3\times 10^{-3}\times 10^{-3}/(\pi\times
5.06\times 10^{-20}) = \qty{1.5e14}{Pa.s/m^3}$; parallel geschakeld
$1.5\times 10^{14}/3\times 10^{6} = \qty{5.0e7}{Pa.s/m^3}$.
\textbf{10.} $8.3\times 10^{-5}\times 5.0\times 10^{7} = \qty{4200}{Pa}
= \qty{31}{mmHg}$.
\textbf{11.} Open \omterm{def:b2:blood-circulation:vessels}{haarvaten} $10^{10}$, elk $\pi(3\times 10^{-6})^{2}
= \qty{2.8e-11}{m^2}$: \qty{0.28}{m^2} $= \qty{2800}{cm^2}$; $v =
83/2800 = \qty{0.03}{cm/s}$; doorlooptijd $0.1/0.03 = \qty{3.3}{s}$.
\textbf{12.} Weerstand $\times(1/0.9)^{4} = 1.52$: \qty{47}{mmHg}; het hart
moet de arteriële druk met \qty{16}{mmHg} verhogen, of het debiet daalt met
een derde.
\textbf{13.} $+10$, $-10$ en $0$ mmHg.
\textbf{14.} Filtratie $10\times 2 = \qty{20}{L}$ per dag; resorptie
$8.5\times 2 = \qty{17}{L}$.
\textbf{15.} \qty{3}{L} lymfe per dag.
\textbf{16.} $\pi_c = \qty{12.5}{mmHg}$: netto $+22.5$ aan het arteriële
uiteinde en $+2.5$ aan het veneuze --- overal filtratie, geen resorptie:
oedeem.
\textbf{17.} Netto $+10$ en $+3$, gemiddeld $+6.5$: filtratie over het hele
\omterm{def:b2:blood-circulation:vessels}{haarvat}, zo'n \qty{26}{L} per dag; de lymfevaten kunnen dat niet afvoeren en
de vloeistof hoopt zich in de weefsels op, eerst in de voeten.
\textbf{18.} $4\times 10^{10}\times 2\pi\times 3\times 10^{-6}\times
10^{-3} = \qty{750}{m^2}$; $t = x^{2}/2D = (2\times 10^{-5})^{2}/
(2\times 10^{-9}) = \qty{0.2}{s}$.
\textbf{19.} $(93 - 3)/83 = \qty{1.08}{mmHg.s/mL}$.
\textbf{20.} $RC = \qty{2.2}{s}$; $120e^{-0.8/2.2} = 120\times 0.69 =
\qty{83}{mmHg}$.
\textbf{21.} $RC = \qty{1.1}{s}$; $120e^{-0.73} = \qty{58}{mmHg}$: de
\omterm{thm:b2:blood-circulation:windkessel}{polsdruk} neemt toe van 37 tot \qty{62}{mmHg}.
\textbf{22.} $120e^{-0.3/2.2} = \qty{105}{mmHg}$: bij een snelle frequentie
daalt de druk tussen de slagen nauwelijks.
\textbf{23.} $C\Delta P = 2\times 20 = \qty{40}{mL}$ opgeslagen; de andere
\qty{30}{mL} stromen al tijdens de systole zelf via de \omterm{def:b2:blood-circulation:vessels}{arteriolen} weg.
\textbf{24.} Het samentrekkende ventrikel drukt in de systole zijn eigen vaten
dicht, zodat de hartspier in de diastole wordt doorbloed, door de druk die de
terugverende aorta in stand houdt; een stijve aorta die de diastolische druk
laat instorten, hongert het hart tussen de slagen uit.
\textbf{25.} Debiet \qty{5}{L/min}; drukval over de \omterm{def:b2:blood-circulation:vessels}{arteriolen} \qty{31}{mmHg};
netto filtraat \qty{3}{L} per dag naar de lymfe; diastolisch \qty{83}{mmHg}
voor $C = 2$ en \qty{58}{mmHg} voor $C = 1$.
\end{solution}
